\section{Small Data and Intelligent Algorithm}

\subsection{Description and Exploratory Data Analysis (EDA)}

The dataset used for this project consists of approximately 225 simulated multi-turn dialogues between a dentist and a patient. Each conversation is annotated with:
\begin{itemize}
    \item the doctor's question,
    \item the corresponding \textit{intent label} (e.g., \texttt{where\_location}, \texttt{duration}, \texttt{trigger}, \texttt{severity}, \texttt{extra}),
    \item and the patient's underlying disease (e.g., \textit{pulpitis acute}, \textit{gingivitis}, \textit{periodontitis}).
\end{itemize}

The dataset was synthetically constructed based on realistic dental symptom descriptions collected from open-access resources such as PubMed and PMC. No personal or identifiable data were used, thus complying with GDPR and ethical data usage regulations.

\subsubsection*{Dataset Summary}
\begin{center}
\begin{tabular}{|l|l|}
\hline
\textbf{Field} & \textbf{Description} \\
\hline
doctor\_question & Natural language question from the dentist \\
turn\_intent & Intent label (type of information requested) \\
disease\_label & Simulated condition of the patient \\
patient\_answer & Generated answer for realism \\
\hline
\end{tabular}
\end{center}

Example records:
\begin{center}
\begin{tabular}{|l|l|l|}
\hline
doctor\_question & turn\_intent & disease \\
\hline
Where exactly does it hurt? & where\_location & pulpitis\_acute \\
Since when do you feel this pain? & duration & pulpitis\_acute \\
What makes the pain worse? & trigger & pulpitis\_acute \\
\hline
\end{tabular}
\end{center}

\noindent
The dataset contains five main intent categories that are approximately balanced. Each question averages between 8 and 12 words, which is ideal for small-scale text classification experiments.

---

\subsection{Intelligent Algorithm Description}

The proposed system is a hybrid AI pipeline composed of two main components:
\begin{enumerate}
    \item A \textbf{Machine Learning intent classifier} trained on the small dataset.
    \item A \textbf{Large Language Model (LLM)} that generates conversational responses as a simulated dental patient.
\end{enumerate}

\subsubsection*{A. Intent Classification Model}

\textbf{Goal:} Identify the intent of the dentist's question (e.g., duration, trigger, severity).

\textbf{Algorithm:} TF-IDF vectorizer + Logistic Regression.

\textbf{Libraries:} \texttt{scikit-learn}, \texttt{pandas}, \texttt{matplotlib}.

\textbf{Motivation:}
\begin{itemize}
    \item Logistic Regression performs well on small datasets and provides interpretable coefficients.
    \item TF-IDF captures key textual patterns such as ``when'' (duration) or ``where'' (location).
    \item The model is simple, efficient, and robust against overfitting.
\end{itemize}

---

\subsubsection*{B. Conversational LLM}

\textbf{Model:} \texttt{mistralai/Mistral-7B-Instruct-v0.2}

\textbf{Framework:} Hugging Face \texttt{transformers}.

\textbf{Quantization:} 4-bit using \texttt{BitsAndBytesConfig} for memory-efficient inference on Colab T4 GPU.

\textbf{Motivation:}
\begin{itemize}
    \item Mistral 7B is instruction-tuned and conversationally fluent.
    \item Generates short, human-like answers without repeating the prompt.
    \item Handles medical dialogue well, simulating realistic patient behavior.
\end{itemize}

---

\subsection{Experimental Methodology and Results}

\subsubsection*{Data Split}
The dataset of 225 examples was divided as follows:
\begin{itemize}
    \item Training set: 80\% 
    \item Validation/Test set: 20\%
\end{itemize}

\subsubsection*{Evaluation Metrics}
For the intent classifier:
\begin{itemize}
    \item Precision, Recall, F1-score per class
    \item Confusion Matrix visualization
\end{itemize}

\subsubsection*{Model Hyperparameters}
\begin{center}
\begin{tabular}{|l|l|}
\hline
Parameter & Value \\
\hline
Vectorizer & TF-IDF \\
C & 1.0 \\
max\_iter & 1000 \\
test\_size & 0.2 \\
random\_state & 42 \\
\hline
\end{tabular}
\end{center}

\subsubsection*{Results}

\begin{verbatim}
Precision: 1
Macro F1-score: 1
Recall: 1
\end{verbatim}

Confusion matrix analysis showed minor overlap between the \texttt{trigger} and \texttt{extra} classes due to semantic similarity.


\subsubsection*{Conversation Generation Results}

After intent prediction, Mistral 7B generates patient responses such as:

\begin{verbatim}
Doctor: Where exactly does it hurt?
Patient (where_location): I'm experiencing a sharp, constant pain in my tooth, especially when 
I bite down or take hot or cold food or drinks. It's been getting worse over the past few days 
and keeps me awake at night.

Doctor: Since when do you feel this pain?
Patient (duration): I've been experiencing intense tooth pain for the past few days. It hurts 
when I apply pressure or eat something cold or hot. The pain keeps me up at night and I can't 
focus on my daily activities due to the discomfort

Doctor: What makes the pain worse?
Patient (trigger): I feel a sharp, constant pain in my tooth when I take a bite or when I 
expose it to cold or hot temperatures. The pain can be quite intense and makes it difficult 
for me to chew on that side of my mouth.

Doctor: How strong is the pain?
Patient (severity): I'm experiencing intense pain in my tooth, especially when I bite down or
take hot or cold food or drinks. The pain keeps me awake at night.

Doctor: Is there anything else you noticed?
Patient (extra): I've been experiencing intense tooth pain, especially when I chew or drink 
something hot or cold. It feels like a sharp, throbbing sensation deep inside my tooth.
Sometimes I even feel a constant ache that keeps me up at night.

\end{verbatim}

---

\subsection{Diagnosis Evaluation Step}

After the conversation, the student (doctor) proposes a diagnosis.  
The LLM knows the true diagnosis from the prompt and compares the student's answer with the ground truth.

\textbf{Example outputs:}
\begin{verbatim}
Correct diagnosis. The symptoms match acute pulpitis.
Incorrect diagnosis. The pain sensitivity suggests pulpitis, not gingivitis.
\end{verbatim}

This step allows students to receive instant feedback on diagnostic accuracy.

---

\subsection{Discussion}

\textbf{Performance summary:}
\begin{itemize}
    \item Mistral 7B produced natural, concise patient responses.
    \item The model runs efficiently in 4-bit mode (~7GB VRAM) on Colab T4 GPU.
\end{itemize}

\textbf{Challenges:}
\begin{itemize}
    \item Smaller models (TinyLlama, Phi-2) often repeated instructions or added meta-text.
    \item Maintaining long-term conversational context remains an open challenge.
\end{itemize}

\textbf{Future Work:}
\begin{itemize}
    \item Add conversation memory through summarization or embeddings.
    \item Include multiple diseases for broader diagnostic practice.
    \item Fine-tune Mistral on real annotated dental dialogues.
\end{itemize}

---

\subsection{Conclusion}

This project demonstrates a hybrid AI pipeline that integrates:
\begin{itemize}
    \item a \textbf{small-data intent classifier} for structured understanding, and
    \item a \textbf{large instruction-tuned LLM (Mistral-7B)} for natural dialogue generation.
\end{itemize}

The resulting system can simulate realistic patient interactions, provide immediate diagnostic feedback, and serve as a valuable training tool for dental students. The combination of explainable ML and generative AI achieves both interpretability and conversational realism within limited computational resources.



\bibliographystyle{plain}
\bibliography{references}






\setstretch{1.5}
\end{document}